\documentclass[a4paper,12pt]{article}
\usepackage[T2A]{fontenc}
\usepackage[utf8]{inputenc}
\usepackage[russian]{babel}
\usepackage{amsmath,amssymb,mathtools}
\usepackage{icomma}
\usepackage[a4paper,top=22mm,bottom=22mm,left=22mm,right=22mm,headheight=15pt]{geometry}
\usepackage{microtype}
\usepackage{tikz}
\usetikzlibrary{arrows.meta,positioning,shapes.geometric,calc}
\usepackage{booktabs,tabularx}
\usepackage{enumitem}
\usepackage{needspace}
\usepackage{fancyhdr}
\usepackage{xcolor}
\usepackage[hidelinks]{hyperref}
\definecolor{accent}{HTML}{008080}
\definecolor{accentdark}{HTML}{006667}
\definecolor{ink}{HTML}{343838}
\definecolor{muted}{HTML}{667070}
\color{ink}
\setlength{\parindent}{0pt}
\setlength{\parskip}{6pt}
\setlength{\headsep}{11pt}
\linespread{1.12}
\pagestyle{fancy}
\fancyhf{}
\fancyhead[L]{\small Финансовые задачи ЕГЭ}
\fancyhead[R]{\small Анастасия Павлова}
\fancyfoot[C]{\small\thepage}
\renewcommand{\headrulewidth}{0.4pt}
\renewcommand{\footrulewidth}{0pt}
\addto\captionsrussian{\renewcommand{\tablename}{Таблица}}
\newcommand{\sect}[1]{\par\needspace{6\baselineskip}\vspace{.8em}\noindent{\large\color{accent}\bfseries #1}\par\vspace{.25em}{\color{accent}\rule{\linewidth}{.35pt}}\par\vspace{.15em}}
\newcommand{\idea}[1]{\textcolor{accentdark}{\textbf{#1}}}
\newcommand{\sample}[1]{\par\textbf{Пример.}\enspace #1\par}
\newcommand{\li}[1]{\par\noindent$\square$\enspace #1}
\setlist[enumerate]{leftmargin=*,itemsep=.75em,topsep=.45em}
\setlist[itemize]{leftmargin=*,itemsep=.28em,topsep=.2em}
\begin{document}

% Обложка — композиция по референсу
\begin{titlepage}
\thispagestyle{empty}
\centering
\vspace*{28mm}
{\large\color{accentdark} МАТЕМАТИКА\par}
\vspace{22mm}
{\LARGE Чек-лист\par}
\vspace{13mm}
{\LARGE\color{accent} Финансовые задачи ЕГЭ:\par}
\vspace{2mm}
{\LARGE\color{accent} кредиты, вклады, проценты\par}
\vspace{21mm}
{\normalsize Подготовка к ЕГЭ по профильной математике\par}
\vfill
{\large Анастасия Павлова\par}
\vspace{13mm}
{\normalsize 5 ноября 2026 г.\par}
\vspace{18mm}
{\small Лёвин Артём Александрович\par}
{\small эксклюзивно для Анастасии Павловой\par}
\vspace{22mm}
\begin{tikzpicture}
\draw[accent,line width=.95pt] (-4.4,0) .. controls (-2.6,.9) and (-1.8,-.6) .. (0,0)
  .. controls (1.7,.7) and (2.9,-.65) .. (4.4,0);
\end{tikzpicture}
\vspace{6mm}
\end{titlepage}

\sect{1. Сначала определите обозначения}

Во всех формулах ниже предполагаются \idea{ежегодное начисление процентов} и платежи или пополнения \idea{в конце года}, если в условии не сказано иначе.

\begin{center}
\begin{tabularx}{\linewidth}{@{}lX@{}}
\toprule
$S$ & исходная сумма кредита или вклада;\\
$n$ & число лет;\\
$r$ & годовая ставка \emph{в процентах};\\
$p=\dfrac{r}{100}$ & та же ставка \emph{в долях единицы};\\
$k=1+p$ & повышающий коэффициент;\\
$x$ & одинаковый ежегодный платёж по кредиту;\\
$A$ & дополнительный взнос на вклад;\\
$F$ & сумма вклада в конце срока;\\
$P_j$ & платёж за $j$-й год при дифференцированной схеме.\\
\bottomrule
\end{tabularx}
\end{center}

\sample{Если ставка $20\%$, то}
\[
 p=\frac{20}{100}=0{,}2,\qquad k=1+0{,}2=1{,}2.
\]

\idea{Важно:} число $1{,}2$ означает увеличение на $20\%$. Число $1{,}22$ соответствует ставке $22\%$.

\sect{2. Аннуитетный кредит: равные платежи}

При аннуитетной схеме \idea{каждый год выплачивается одна и та же сумма $x$}. Сначала начисляются проценты на остаток долга, затем вносится платёж.

Пусть $B_j$ --- остаток долга после $j$-го платежа. Тогда
\[
 B_0=S,\qquad B_j=kB_{j-1}-x,\qquad B_n=0.
\]

Последовательно раскрывая эту запись, получаем основную формулу:
\[
 \boxed{S k^n=x\bigl(k^{n-1}+k^{n-2}+\ldots+k+1\bigr).}
\]

При двух и трёх годах соответственно:
\[
 \begin{aligned}
 n=2:\quad &Sk^2-xk-x=0;\\[2pt]
 n=3:\quad &Sk^3-xk^2-xk-x=0.
 \end{aligned}
\]

\idea{Не перепутайте степень:} срок в три года означает $k^3$ и три слагаемых с платежом $x$.

\sect{3. Как найти платёж и переплату}

\sample{Кредит $182\,000$ руб. выдан на три года под $20\%$ годовых. Ежегодный платёж одинаков. Найдите платёж, общую выплату и переплату.}

Сначала записываем формулу в общем виде:
\[
 Sk^3=x(k^2+k+1).
\]
Подставляем $S=182\,000$, $k=1{,}2$:
\[
 182\,000\cdot1{,}2^3=x(1{,}2^2+1{,}2+1),
\]
\[
 314\,496=3{,}64x,\qquad x=86\,400\text{ руб.}
\]

\[
 3x=259\,200\text{ руб.},\qquad 3x-S=77\,200\text{ руб.}
\]

\idea{Проверка:} после третьего платежа остаток долга равен нулю.

Для произвольного срока обязательно различайте три величины:
\[
 \boxed{\text{платёж }x}\qquad
 \boxed{\text{все платежи }nx}\qquad
 \boxed{\text{переплата }nx-S}.
\]

\sect{4. Вклад с пополнением после первого года}

Если исходный вклад $S$ размещён на два года, а \idea{в конце первого года} внесено дополнительно $A$, то к окончанию второго года:
\[
 \boxed{F=Sk^2+Ak.}
\]

Первоначальная сумма работает два года, поэтому получает множитель $k^2$. Дополнительный взнос работает только один год, поэтому получает множитель $k$.

\sample{Вложили $500\,000$ руб., после первого года добавили $200\,000$ руб. Через два года получили $851\,200$ руб. Найдите годовую ставку.}
\[
 500\,000k^2+200\,000k=851\,200,
\]
Для определения $k$ решаем квадратное уравнение:
\[
\begin{aligned}
1250k^2+500k-2128&=0,\\
D=500^2+4\cdot1250\cdot2128&=3300^2,\\
k=\frac{-500\pm3300}{2500}&\in\{1{,}12;\,-1{,}52\}.
\end{aligned}
\]
Отрицательный корень не подходит. При $k=1{,}12$ проверка даёт
$1250\cdot1{,}12^2+500\cdot1{,}12=2128$.
Значит, $p=k-1=0{,}12$, откуда \idea{$r=12\%$}.

При решении квадратного уравнения выбирайте корень, соответствующий условию: для положительной процентной ставки $k>1$. Если ставка может быть нулевой, $k\ge1$.

\sect{5. Дифференцированный кредит: равные доли долга}

При этой схеме \idea{основной долг ежегодно уменьшается на одну и ту же сумму $S/n$}. Проценты начисляют на долг, оставшийся \emph{перед} очередным платежом.

\[
 \boxed{P_j=\frac Sn +p\left(S-\frac{(j-1)S}{n}\right).}
\]

Для первого и последнего платежей удобно запомнить:
\[
 \boxed{P_1=\frac Sn+Sp,\qquad P_n=\frac Sn+\frac Sn\,p.}
\]

Разница соседних платежей постоянна:
\[
 P_j-P_{j+1}=\frac Sn\,p.
\]

Полная выплата по кредиту:
\[
 \boxed{\sum_{j=1}^{n}P_j=S+Sp\frac{n+1}{2}.}
\]

\sample{Кредит $1\,200\,000$ руб. на четыре года при ставке $7\%$.}

Каждый год погашается $1\,200\,000/4=300\,000$ руб. основного долга. Остатки перед начислением процентов: $1\,200\,000$; $900\,000$; $600\,000$; $300\,000$ руб.

\[
 \begin{aligned}
 P_1&=300\,000+84\,000=384\,000,\\
 P_2&=300\,000+63\,000=363\,000,\\
 P_3&=300\,000+42\,000=342\,000,\\
 P_4&=300\,000+21\,000=321\,000.
 \end{aligned}
\]

Все платежи: $1\,410\,000$ руб.; переплата: $210\,000$ руб.

\sect{6. Как работать с ограничениями на платежи}

\sample{Кредит $1\,200\,000$ руб. на четыре года. Первый платёж не больше $390\,000$ руб., последний --- не меньше $320\,000$ руб. Ставка в процентах является целым числом. Найдите её.}

Записываем общие формулы \emph{до} подстановки чисел:
\[
 \begin{cases}
 P_1=\dfrac{S}{n}+Sp\le390\,000,\\[5pt]
 P_4=\dfrac{S}{n}+\dfrac{S}{n}p\ge320\,000.
 \end{cases}
\]

После подстановки $S=1\,200\,000$, $n=4$:
\[
 \begin{cases}
 300\,000+1\,200\,000p\le390\,000,\\
 300\,000+300\,000p\ge320\,000.
 \end{cases}
\]
Отсюда
\[
 \frac1{15}\le p\le\frac3{40}.
\]

Переходим от долей единицы к процентам (умножаем на $100$):
\[
 6\frac23\%\le r\le7{,}5\%.
\]

Единственная целая ставка в этом промежутке --- \idea{$7\%$}. Проверка: $P_1=384\,000\le390\,000$; $P_4=321\,000\ge320\,000$.

\sect{7. Типичные ошибки и самопроверка}
\li{Проверили число лет? При $n=3$ уравнение аннуитета содержит три платежа.}
\li{Процент перевели в долю единицы? Например, $25\%=0{,}25$ и $k=1{,}25$.}
\li{Нашли именно то, что спрашивают: платёж, полную сумму или переплату?}
\li{Для вклада правильно указали момент дополнительного взноса?}
\li{В дифференцированном кредите проценты считали от текущего остатка долга?}
\li{Если требуется целая ставка, нашли все целые числа из допустимого интервала?}
\li{Проверили ответ подстановкой и указали рубли или проценты?}

\clearpage
\thispagestyle{plain}
\begin{center}
{\LARGE\color{accent} Алгоритм: как распознать финансовую схему}\par
\vspace{5mm}
Условие задачи определяет, какую формулу нужно составить.
\end{center}
\vspace{12mm}
\begin{center}
\begin{tikzpicture}[
 >=Stealth,
 start/.style={ellipse,draw=accentdark,thick,minimum width=4.9cm,minimum height=1.05cm,align=center,font=\small},
 decide/.style={diamond,aspect=2.15,draw=accent,thick,text width=3.2cm,inner sep=3pt,align=center,font=\small},
 action/.style={rectangle,rounded corners=2pt,draw=accent,thick,text width=5.65cm,minimum height=1.35cm,inner sep=8pt,align=center,font=\small},
 next/.style={rectangle,rounded corners=2pt,draw=accentdark,thick,text width=4.85cm,minimum height=1.3cm,inner sep=8pt,align=center,font=\small},
 arr/.style={-{Stealth[length=2.6mm]},line width=.85pt,color=accentdark},
 lab/.style={font=\small,fill=white,inner sep=2pt}
]
\node[start] (begin) at (0,0) {Прочитайте условие и цель задачи};
\node[decide] (d1) at (0,-2.65) {Деньги размещены на вкладе?};
\node[action] (a1) at (6.1,-2.65) {\textbf{Вклад}\quad $F=Sk^2+Ak$\\[2pt] Уточните момент пополнения};
\node[decide] (d2) at (0,-7.05) {Ежегодные платежи одинаковы?};
\node[action] (a2) at (6.1,-7.05) {\textbf{Аннуитет}\quad $Sk^n=x\sum_{j=0}^{n-1}k^j$\\[2pt] Затем найдите $nx$ и $nx-S$};
\node[decide] (d3) at (0,-11.45) {Основной долг уменьшается равными долями?};
\node[action] (a3) at (6.1,-11.45) {\textbf{Дифференцированная схема}\\[2pt] $P_j=\dfrac Sn+p\left(S-\dfrac{(j-1)S}{n}\right)$};
\node[next] (other) at (0,-15.7) {Другая схема: сначала опишите, как изменяется остаток долга};
\draw[arr] (begin.south)--(d1.north);
\draw[arr] (d1.east)--node[lab,above]{да}(a1.west);
\draw[arr] (d1.south)--node[lab,right]{нет}(d2.north);
\draw[arr] (d2.east)--node[lab,above]{да}(a2.west);
\draw[arr] (d2.south)--node[lab,right]{нет}(d3.north);
\draw[arr] (d3.east)--node[lab,above]{да}(a3.west);
\draw[arr] (d3.south)--node[lab,right]{нет}(other.north);
\end{tikzpicture}
\end{center}
\clearpage

\begin{center}
{\LARGE\color{accent} Тренировочный блок}\par
\vspace{1mm}
{\normalsize Путь А --- 10 заданий, от базового к повышенному уровню}\par
\end{center}

Во всех задачах считайте, что проценты начисляются раз в год, а платежи и дополнительные взносы производятся в конце года. Если требуется, указывайте денежные суммы в рублях.

\begin{enumerate}
\item Годовая ставка по кредиту составляет $18\%$. Найдите $p$ (долю единицы) и повышающий коэффициент $k$.

\item Определите схему для каждого описания: а) все годовые платежи по кредиту равны; б) ежегодно погашается одинаковая часть основного долга; в) остаток основного долга каждый год сокращается на $25\%$ от текущего остатка. Для пункта в) поясните, относится ли условие само по себе к одной из двух стандартных схем.

\item Кредит $165\,000$ руб. выдан на два года под $20\%$ годовых с двумя равными ежегодными платежами. Найдите ежегодный платёж, всю сумму платежей и переплату.

\item Кредит $91\,000$ руб. выдан на три года под $20\%$ годовых. Платежи в конце каждого года равны. Найдите ежегодный платёж, всю выплату и переплату.

\item На вклад внесли $400\,000$ руб. под $10\%$ годовых. В конце первого года внесли ещё $150\,000$ руб. Сколько денег будет на вкладе в конце второго года?

\item На вклад внесли $300\,000$ руб. В конце первого года добавили $100\,000$ руб. После второго года на счёте стало $473\,000$ руб. Найдите годовую процентную ставку, если она положительна и одинакова в оба года.

\item Кредит $900\,000$ руб. взят на три года под $6\%$ годовых с равным ежегодным погашением основного долга. Найдите все три годовых платежа, общую выплату и переплату.

\item Кредит $800\,000$ руб. взят на четыре года под $5\%$ годовых. Основной долг погашается равными частями. На сколько рублей каждый следующий ежегодный платёж меньше предыдущего?

\item Кредит $1\,000\,000$ руб. взят на пять лет с равным ежегодным погашением основного долга. Первый платёж не превышает $277\,000$ руб., а последний составляет не менее $214\,000$ руб. Ставка в процентах --- положительное целое число. Найдите ставку и общую сумму выплат.

\item Кредит $180\,000$ руб. погашается двумя равными ежегодными платежами по $125\,000$ руб. Найдите положительную годовую процентную ставку.
\end{enumerate}

\clearpage
\begin{center}
{\LARGE\color{accent} Ответы}\par
\vspace{2mm}
Таблица к тренировочному блоку «Путь А»
\end{center}
\vspace{5mm}
\renewcommand{\arraystretch}{1.45}
\begin{tabularx}{\linewidth}{@{}cX@{}}
\toprule
\textbf{№} & \textbf{Ответ}\\
\midrule
1 & $p=0{,}18$; $k=1{,}18$.\\
2 & а) аннуитет; б) дифференцированная схема; в) отдельное правило изменения остатка долга, которого недостаточно для определения стандартной схемы.\\
3 & Платёж $108\,000$ руб.; все платежи $216\,000$ руб.; переплата $51\,000$ руб.\\
4 & Платёж $43\,200$ руб.; все платежи $129\,600$ руб.; переплата $38\,600$ руб.\\
5 & $649\,000$ руб.\\
6 & $10\%$.\\
7 & $354\,000$; $336\,000$; $318\,000$ руб.; всего $1\,008\,000$ руб.; переплата $108\,000$ руб.\\
8 & На $10\,000$ руб.\\
9 & Ставка $7\%$; всего $1\,210\,000$ руб.\\
10 & $25\%$.\\
\bottomrule
\end{tabularx}
\vspace{16mm}
\begin{center}
\textcolor{accentdark}{\large Итоговая памятка}
\end{center}
\begin{tabularx}{\linewidth}{@{}p{.34\linewidth}X@{}}
\toprule
\textbf{Сигнал в условии} & \textbf{Действие}\\
\midrule
«Равные ежегодные платежи» & Используйте формулу аннуитета $Sk^n=x\sum_{j=0}^{n-1}k^j$.\\
«Равное погашение основного долга» & Добавьте к $S/n$ проценты на текущий остаток долга.\\
«Пополнение вклада» & Учитывайте, сколько лет начисляются проценты на каждый взнос.\\
«Переплата» & Вычтите исходную сумму кредита из всех платежей.\\
«Целая ставка» & Составьте неравенства и отберите целые проценты из интервала.\\
\bottomrule
\end{tabularx}

\end{document}
